\setlength{\englishIndent}{0pt}

\chapter{Añjali}

Il canto e le richieste formali si eseguono con le mani in añjali.
Questo è un gesto di rispetto, che si compie unendo i palmi delle mani
direttamente davanti al petto, con le dita allineate e rivolte
verso l'alto.

\chapter{Richiesta di un discorso di Dhamma}

\begin{instruction}
  Dopo essersi inchinati tre volte, con le mani giunte in añjali,\\
  recitare quanto segue:
\end{instruction}

Brahmā ca꜕ lokādhipa꜕tī sa꜕hampa꜕ti\\
Ka꜕tañja꜕lī a꜕nadhiva꜕raṁ a꜕yāca꜕tha

Santī꜓dha sa꜕ttāppa꜕ra꜕jakkha꜕-jātikā\\
Desetu꜕ dhammaṁ a꜕nu꜕kampi꜕maṁ pa꜕jaṁ

\begin{instruction}
  Inchinatevi di nuovo tre volte
\end{instruction}

\begin{english}
Il dio Brahmā Sahampati, Signore del mondo,\\
con le mani giunte in segno di riverenza, chiese un favore:

«Ci sono esseri qui con poca polvere negli occhi,\\
insegnate il Dhamma per compassione verso di loro».
\end{english}

\chapter{Riconoscimento dell'insegnamento}

\enlargethispage{\baselineskip}

\begin{tabular}{@{} ll @{}}
Una persona: & Ha꜓nda mayaṁ dhammakathā꜓ya sā꜓dhukā꜕raṁ dadāmase \\
& \hspace*{1em}\tr{Ora esprimiamo la nostra approvazione per questo insegnamento di Dhamma.} \\
Risposta: & Sādhu, sādhu, sādhu, anu꜓modāmi \\
& \hspace*{1em}\tr{Così sia, lo apprezzo.} \\
\end{tabular}

\clearpage
\chapter{Richiesta dei canti Paritta}

\begin{instruction}
  Dopo essersi inchinati tre volte, con le mani giunte in añjali,\\
  recitare quanto segue:
\end{instruction}

Vipatti-paṭibāhā꜓ya sabba꜕-sampatti꜕-siddhi꜕yā\\
Sabbadukkha-vināsā꜓ya\\
Parittaṁ brūtha꜕ maṅga꜕laṁ

Vipatti-paṭibāhā꜓ya sabba꜕-sampatti꜕-siddhi꜕yā\\
Sabbabhaya-vināsā꜓ya\\
Parittaṁ brūtha꜕ maṅga꜕laṁ

Vipatti-paṭibāhā꜓ya sabba꜕-sampatti꜕-siddhi꜕yā\\
Sabbaroga-vināsā꜓ya\\
Parittaṁ brūtha꜕ maṅga꜕laṁ

\begin{instruction}
  Inchinatevi tre volte
\end{instruction}

\begin{english}
Per allontanare la sfortuna, per il sorgere della buona sorte,\\
per la dissoluzione di ogni dukkha,\\
possiate cantare una benedizione e protezione.

Per allontanare la sfortuna, per il sorgere della buona sorte,\\
per la dissoluzione di ogni paura,\\+possiate cantare una benedizione e protezione.

Per allontanare la sfortuna, per il sorgere della buona sorte,\\
per la dissoluzione di ogni malattia,\\
possiate cantare una benedizione e protezione.
\end{english}

\setlength{\englishIndent}{\leaderIndent}

\clearpage
\chapter[Tre Rifugi e Cinque Precetti]{Richiesta dei Tre Rifugi\newline e dei Cinque Precetti}

\begin{instruction}
  Dopo essersi inchinati tre volte, con le mani giunte in añjali,\\
  recitare la richiesta appropriata.
\end{instruction}

\section{Per un gruppo, a un monaco}

\begin{twochants}
Mayaṁ bhante tisaraṇena sa꜕ha & pañca sī꜓lāni yā꜕cāma\\
Dutiyampi mayaṁ bhante tisaraṇena sa꜕ha & pañca sī꜓lāni yā꜕cāma\\
Tatiyampi mayaṁ bhante tisaraṇena sa꜕ha & pañca sī꜓lāni yā꜕cāma\\
\end{twochants}

\section{Per se stessi, a un monaco}

\begin{twochants}
Ahaṁ bhante tisaraṇena sa꜕ha & pañca sī꜓lāni yā꜕cāmi\\
Dutiyampi ahaṁ bhante tisaraṇena sa꜕ha & pañca sī꜓lāni yā꜕cāmi\\
Tatiyampi ahaṁ bhante tisaraṇena sa꜕ha & pañca sī꜓lāni yā꜕cāmi
\end{twochants}

\section{Per un gruppo, a una monaca}

\begin{twochants}
Mayaṁ ayye tisaraṇena sa꜕ha & pañca sī꜓lāni yā꜕cāma\\
Dutiyampi mayaṁ ayye tisaraṇena sa꜕ha & pañca sī꜓lāni yā꜕cāma\\
Tatiyampi mayaṁ ayye tisaraṇena sa꜕ha & pañca sī꜓lāni yā꜕cāma\\
\end{twochants}

\section{Per se stessi, a una monaca}

\begin{twochants}
Ahaṁ ayye tisaraṇena sa꜕ha & pañca sī꜓lāni yā꜕cāmi\\
Dutiyampi ahaṁ ayye tisaraṇena sa꜕ha & pañca sī꜓lāni yā꜕cāmi\\
Tatiyampi ahaṁ ayye tisaraṇena sa꜕ha & pañca sī꜓lāni yā꜕cāmi\\
\end{twochants}

\section{Per un gruppo, a un laico/a}

\begin{twochants}
Mayaṁ mitta tisaraṇena sa꜕ha & pañca sī꜓lāni yā꜕cāma\\
Dutiyampi mayaṁ mitta tisaraṇena sa꜕ha & pañca sī꜓lāni yā꜕cāma\\
Tatiyampi mayaṁ mitta tisaraṇena sa꜕ha & pañca sī꜓lāni yā꜕cāma\\
\end{twochants}

\section{Per se stessi, a un laico/a}

\begin{twochants}
Ahaṁ mitta tisaraṇena sa꜕ha & pañca sī꜓lāni yā꜕cāmi\\
Dutiyampi ahaṁ mitta tisaraṇena sa꜕ha & pañca sī꜓lāni yā꜕cāmi\\
Tatiyampi ahaṁ mitta tisaraṇena sa꜕ha & pañca sī꜓lāni yā꜕cāmi\\
\end{twochants}

\section{Traduzione}

\begin{english}
  Noi/Io, Venerabile Signore/Signora/Amico/a,\\
  richiediamo i Tre Rifugi e i Cinque Precetti.

  Per la seconda volta,\\
  noi/io, Venerabile Signore/Signora/Amico/a,\\
  richiediamo i Tre Rifugi e i Cinque Precetti.

  Per la terza volta,\\
  noi/io, Venerabile Signore/Signora/Amico/a,\\
  richiediamo i Tre Rifugi e i Cinque Precetti.
\end{english}

\clearpage
\chapter{Prendere i Tre Rifugi}

\begin{instruction}
  Ripetere, dopo che il leader ha cantato le prime tre righe
\end{instruction}

Namo tassa bhagavato arahato sammāsambuddhassa\\
Namo tassa bhagavato arahato sammāsambuddhassa\\
Namo tassa bhagavato arahato sammāsambuddhassa

\begin{english}
  Omaggio al Beato, Nobile e Perfettamente Risvegliato.\\
  Omaggio al Beato, Nobile e Perfettamente Risvegliato.\\
  Omaggio al Beato, Nobile e Perfettamente Risvegliato.
\end{english}

Buddhaṁ saraṇaṁ gacchāmi\\
Dhammaṁ saraṇaṁ gacchāmi\\
Saṅghaṁ saraṇaṁ gacchāmi

\begin{english}
  Prendo rifugio nel Buddha.\\
  Prendo rifugio nel Dhamma.\\
  Prendo rifugio nel Saṅgha.
\end{english}

Dutiyampi buddhaṁ saraṇaṁ gacchāmi\\
Dutiyampi dhammaṁ saraṇaṁ gacchāmi\\
Dutiyampi saṅghaṁ saraṇaṁ gacchāmi

\begin{english}
  Per la seconda volta, prendo rifugio nel Buddha.\\
  Per la seconda volta, prendo rifugio nel Dhamma.\\
  Per la seconda volta, prendo rifugio nel Saṅgha.
\end{english}

Tatiyampi buddhaṁ saraṇaṁ gacchāmi\\
Tatiyampi dhammaṁ saraṇaṁ gacchāmi\\
Tatiyampi saṅghaṁ saraṇaṁ gacchāmi

\clearpage

\begin{english}
  Per la terza volta, prendo rifugio nel Buddha.\\
  Per la terza volta, prendo rifugio nel Dhamma.\\
  Per la terza volta, prendo rifugio nel Saṅgha.
\end{english}

\begin{instruction}
  Leader:
\end{instruction}

[Tisaraṇa-gamanaṁ niṭṭhitaṁ]

\begin{english}
  Questo completa la presa dei Tre Rifugi.
\end{english}

\begin{instruction}
  Risposta:
\end{instruction}

Āma bhante / Āma ayye / Āma mitta

\begin{english}
  Sì, Venerabile Signore/Signora/Amico/a.
\end{english}

\chapter{I Cinque Precetti}

\begin{instruction}
  Ripetere ogni precetto dopo il leader
\end{instruction}

\begin{precept}
  \setcounter{enumi}{0}
  \item Pāṇātipātā vera꜓maṇī sikkhā꜓padaṁ sa꜓mādi꜕yāmi
\end{precept}

\begin{english}
  Mi impegno a osservare il precetto di astenermi dal togliere la vita a qualsiasi creatura vivente.
\end{english}

\begin{precept}
  \setcounter{enumi}{1}
  \item Adinnādānā vera꜓maṇī sikkhā꜓padaṁ sa꜓mādi꜕yāmi
\end{precept}

\begin{english}
  Mi impegno a osservare il precetto di astenermi dal prendere ciò che non è dato.
\end{english}

\begin{precept}
  \setcounter{enumi}{2}
  \item Kāmesu micchā꜓cārā vera꜓maṇī sikkhā꜓padaṁ sa꜓mādi꜕yāmi
\end{precept}

\begin{english}
  Mi impegno a osservare il precetto di astenermi da una condotta sessuale scorretta.
\end{english}

\clearpage

\begin{precept}
  \setcounter{enumi}{3}
  \item Musā꜓vādā vera꜓maṇī sikkhā꜓padaṁ sa꜓mādi꜕yāmi
\end{precept}

\begin{english}
  Mi impegno a osservare il precetto di astenermi dal mentire.
\end{english}

\begin{precept}
  \setcounter{enumi}{4}
  \item Surāmeraya-majja-pamādaṭṭhā꜓nā vera꜓maṇī sikkhā꜓padaṁ sa꜓mādi꜕yāmi
\end{precept}

\begin{english}
  Mi impegno a osservare il precetto di astenermi dal consumare bevande inebrianti e droghe che portano alla negligenza.
\end{english}

\begin{instruction}
  Responsabile:
\end{instruction}

[Imāni pañca sikkhā꜓padāni\\
Sī꜓lena suga꜕tiṁ yanti\\
Sī꜓lena bhoga꜕sa꜓mpadā\\
Sī꜓lena nibbu꜕tiṁ yanti\\
Tasmā꜓ sī꜓laṁ viso꜓dhaye]

\begin{english}
  Questi sono i Cinque Precetti;\\
  la virtù è fonte di felicità,\\
  la virtù è fonte di vera ricchezza,\\
  la virtù è fonte di pace ---\\
  Perciò la virtù sia purificata.
\end{english}

\begin{instruction}
  Risposta:
\end{instruction}

Sādhu, sādhu, sādhu

\begin{instruction}
  Inchinatevi tre volte
\end{instruction}
